\documentclass[a4paper,12pt]{article}

\usepackage[T2A]{fontenc}
\usepackage[utf8]{inputenc}
\usepackage[russian]{babel}
\usepackage{amsmath,amssymb,amsthm,mathtools}
\usepackage{helvet}
\renewcommand{\familydefault}{\sfdefault}
\usepackage{icomma}
\usepackage{geometry}
\usepackage{microtype}
\usepackage{tikz}
\usepackage{tkz-euclide}
\usepackage{pgfplots}
\pgfplotsset{compat=1.18}
\usepackage{tabularx,booktabs,longtable}
\usepackage{enumitem}
\usepackage{hyperref}
\usepackage{parskip}
\usepackage{fancyhdr}
\usepackage{setspace}
\usepackage{xcolor}

\geometry{margin=2.1cm}
\onehalfspacing

% ------------------------------------------------------------
% Палитра
% ------------------------------------------------------------

\definecolor{accent}{HTML}{008080}
\definecolor{darkaccent}{HTML}{004D4D}
\definecolor{textgray}{HTML}{333333}
\definecolor{softgray}{HTML}{777777}
\definecolor{lightgray}{HTML}{E8EEEE}
\definecolor{warmgray}{HTML}{F4F3F0}

\color{textgray}

% ------------------------------------------------------------
% Настройка документа
% ------------------------------------------------------------

\setlength{\parindent}{0pt}
\setlength{\parskip}{0.55em}

\setlist[itemize]{
    leftmargin=1.45em,
    itemsep=0.25em,
    topsep=0.25em
}

\setlist[enumerate]{
    leftmargin=1.7em,
    itemsep=0.45em,
    topsep=0.35em
}

\hypersetup{
    colorlinks=true,
    linkcolor=darkaccent,
    urlcolor=accent,
    pdfauthor={Лёвин Артём Александрович},
    pdftitle={Чек-лист: комплексное повторение алгебры ОГЭ}
}

\pagestyle{fancy}
\fancyhf{}
\fancyhead[L]{\small\textcolor{softgray}{София Кальней}}
\fancyhead[R]{\small\textcolor{softgray}{ОГЭ по математике}}
\fancyfoot[C]{\small\textcolor{softgray}{\thepage}}
\renewcommand{\headrulewidth}{0.35pt}
\renewcommand{\headrule}{\hbox to\headwidth{\color{lightgray}\leaders\hrule height \headrulewidth\hfill}}

\newcommand{\topic}[1]{%
    \vspace{0.5em}
    {\Large\bfseries\textcolor{darkaccent}{#1}}\par
    \vspace{0.2em}
    {\color{accent}\rule{\linewidth}{0.7pt}}
    \vspace{0.5em}
}

\newcommand{\subtopic}[1]{%
    \vspace{0.5em}
    {\large\bfseries\textcolor{darkaccent}{#1}}\par
}

\newcommand{\keyidea}[1]{%
    \par\smallskip
    \textbf{\textcolor{darkaccent}{Ключевая идея.}} #1
    \par\smallskip
}

\newcommand{\warning}[1]{%
    \par\smallskip
    \textbf{\textcolor{accent}{Важно.}} #1
    \par\smallskip
}

\newcommand{\algstep}[1]{%
    \textbf{\textcolor{darkaccent}{#1}}
}

\newcommand{\R}{\mathbb{R}}

% ------------------------------------------------------------
% Документ
% ------------------------------------------------------------

\begin{document}

% ============================================================
% Титульный лист
% ============================================================

\thispagestyle{empty}

\begin{center}

\vspace*{2.5cm}

{\Large\textcolor{accent}{\textbf{ЧЕК-ЛИСТ}}}

\vspace{0.8cm}

{\huge\bfseries\textcolor{darkaccent}{
Комплексное повторение\\[0.15em]
алгебры ОГЭ
}}

\vspace{0.8cm}

{\large
Уравнения, системы, неравенства,\\
текстовые задачи и графики
}

\vspace{1.7cm}

{\large
\textbf{Лёвин Артём Александрович}\\
эксклюзивно для Софии Кальней
}

\vfill

{\large \today}

\vspace{1.4cm}

\begin{tikzpicture}[x=1cm,y=1cm]
    \draw[
        accent,
        line width=1.15pt
    ]
    (-7,0)
    .. controls (-4.8,0.6) and (-2.6,-0.5) ..
    (0,0)
    .. controls (2.7,0.55) and (4.9,-0.45) ..
    (7,0);
\end{tikzpicture}

\vspace{0.8cm}

\end{center}

\newpage

% ============================================================
% Введение
% ============================================================

\topic{1. Карта занятия}

На занятии повторяются основные типы алгебраических задач повышенной сложности ОГЭ. Главная цель --- научиться быстро распознавать структуру задачи и выбирать короткий, надёжный метод решения.

\begin{itemize}
    \item разложение выражений на множители;
    \item группировка и двойное вынесение;
    \item уравнения вида \(A(x)^2=c\);
    \item замена переменной;
    \item системы уравнений;
    \item метод интервалов;
    \item дробно-рациональные неравенства;
    \item текстовые задачи на движение;
    \item задачи на производительность;
    \item смеси, растворы и высушивание;
    \item графики рациональных функций после сокращения;
    \item выколотые точки и задачи с параметром.
\end{itemize}

\keyidea{
На экзамене важно увидеть математическую конструкцию раньше, чем начинать длинные вычисления.
}

% ============================================================
% Факторизация
% ============================================================

\topic{2. Группировка и двойное вынесение}

Если выражение состоит из четырёх слагаемых, полезно проверить возможность группировки по два слагаемых.

Типовая схема:

\[
ab+ac-db-dc.
\]

Группируем:

\[
(ab+ac)-(db+dc).
\]

Во второй группе минус относится ко всей скобке:

\[
a(b+c)-d(b+c).
\]

Теперь появляется общий множитель:

\[
(b+c)(a-d).
\]

\keyidea{
Четыре слагаемых часто подсказывают приём
\[
\text{группировка}\;\longrightarrow\;
\text{вынесение}\;\longrightarrow\;
\text{повторное вынесение}.
\]
}

\warning{
Если перед группой стоит минус, он должен быть учтён для каждого слагаемого внутри скобки.
}

\subtopic{Пример}

Разложим на множители:

\[
x^2+x-4x-4.
\]

Группируем:

\[
(x^2+x)-(4x+4).
\]

Выносим общие множители:

\[
x(x+1)-4(x+1).
\]

Следовательно,

\[
\boxed{(x+1)(x-4)}.
\]

% ============================================================
% Уравнения A^2 = c
% ============================================================

\topic{3. Уравнения со скобкой в квадрате}

Рассмотрим уравнение

\[
(x-4)^2-9=0.
\]

Можно раскрыть квадрат и перейти к квадратному уравнению. Быстрее использовать структуру исходного выражения:

\[
(x-4)^2=9.
\]

Отсюда

\[
x-4=3
\]

или

\[
x-4=-3.
\]

Получаем

\[
x=7,\qquad x=1.
\]

\subtopic{Общая схема}

Если

\[
A(x)^2=c,
\]

то возможны три ситуации.

\begin{enumerate}
    \item \(c>0\):
    \[
    A(x)=\sqrt c
    \quad\text{или}\quad
    A(x)=-\sqrt c.
    \]

    \item \(c=0\):
    \[
    A(x)=0.
    \]

    \item \(c<0\):
    \[
    \boxed{\text{решений в действительных числах нет}}.
    \]
\end{enumerate}

\subtopic{Пример}

\[
(x-5)^2-10=0.
\]

Тогда

\[
(x-5)^2=10,
\]

поэтому

\[
x-5=\pm\sqrt{10}.
\]

Ответ:

\[
\boxed{x=5-\sqrt{10};\quad x=5+\sqrt{10}}.
\]

\warning{
Иррациональные корни вида \(5\pm\sqrt{10}\) являются обычным корректным ответом.
}

\subtopic{Как мгновенно увидеть отсутствие решений}

Рассмотрим

\[
(x-5)^2+10=0.
\]

При любом действительном \(x\)

\[
(x-5)^2\ge 0,
\]

следовательно,

\[
(x-5)^2+10\ge 10>0.
\]

Левая часть не может стать равной нулю:

\[
\boxed{\emptyset}.
\]

% ============================================================
% Замена переменной
% ============================================================

\topic{4. Замена переменной}

Рассмотрим уравнение

\[
\frac{1}{(x-1)^2}
-\frac{5}{x-1}
+6=0.
\]

Сначала фиксируем ограничение:

\[
x\ne1.
\]

В выражении повторяется элемент

\[
\frac{1}{x-1}.
\]

Введём замену

\[
t=\frac{1}{x-1}.
\]

Тогда

\[
t^2-5t+6=0.
\]

Разложим на множители:

\[
(t-2)(t-3)=0.
\]

Следовательно,

\[
t=2
\quad\text{или}\quad
t=3.
\]

Возвращаемся к \(x\):

\[
\frac{1}{x-1}=2,
\qquad
\frac{1}{x-1}=3.
\]

Получаем

\[
x=\frac32,
\qquad
x=\frac43.
\]

\keyidea{
Сильный сигнал для замены --- один и тот же сложный элемент встречается несколько раз и в разных степенях.
}

\subtopic{Биквадратная структура}

Если встречается выражение

\[
(x-4)^4-4(x-4)^2+3=0,
\]

полезно ввести

\[
t=(x-4)^2.
\]

Так как квадрат неотрицателен,

\[
t\ge0.
\]

Получаем

\[
t^2-4t+3=0,
\]

откуда

\[
t=1
\quad\text{или}\quad
t=3.
\]

Далее решаются уравнения

\[
(x-4)^2=1,
\qquad
(x-4)^2=3.
\]

\warning{
После замены вида \(t=A(x)^2\) обязательно учитывайте условие
\[
t\ge0.
\]
Отрицательное значение \(t\) не даёт действительных решений.
}

% ============================================================
% Системы
% ============================================================

\topic{5. Системы уравнений: разветвление и подстановка}

Особенно удобны системы, где одно уравнение уже разложено на множители.

Например:

\[
\begin{cases}
(x-3)y=0,\\[0.3em]
\dfrac{y+1}{x+y-4}=2.
\end{cases}
\]

Сначала учитываем ОДЗ:

\[
x+y-4\ne0.
\]

Первое уравнение обнуляется в двух случаях:

\[
x-3=0
\]

или

\[
y=0.
\]

Следовательно, рассматриваем две ветви:

\[
x=3
\qquad\text{или}\qquad
y=0.
\]

\subtopic{Ветвь 1: \(x=3\)}

Подставляем во второе уравнение:

\[
\frac{y+1}{y-1}=2.
\]

Условие:

\[
y\ne1.
\]

Решаем:

\[
y+1=2y-2,
\]

\[
y=3.
\]

Получаем пару

\[
(3;3).
\]

\subtopic{Ветвь 2: \(y=0\)}

Подставляем:

\[
\frac{1}{x-4}=2.
\]

Условие:

\[
x\ne4.
\]

Тогда

\[
1=2x-8,
\]

\[
x=\frac92.
\]

Вторая пара:

\[
\left(\frac92;0\right).
\]

Ответ:

\[
\boxed{
(3;3),\;
\left(\frac92;0\right)
}.
\]

\keyidea{
В системе сначала удобно использовать уравнение, которое быстрее всего ограничивает возможные значения переменных.
}

\warning{
Ответ системы --- набор упорядоченных пар \((x;y)\). Каждая найденная пара должна удовлетворять всем исходным уравнениям и ОДЗ.
}

% ============================================================
% Метод интервалов
% ============================================================

\topic{6. Метод интервалов}

Рассмотрим неравенство

\[
(x-3)(x^2+x-6)\le0.
\]

Разложим квадратный трёхчлен:

\[
x^2+x-6=(x+3)(x-2).
\]

Получаем

\[
(x-3)(x+3)(x-2)\le0.
\]

Критические точки:

\[
-3,\qquad 2,\qquad 3.
\]

Все множители линейные и входят в первой степени, поэтому знак выражения меняется при прохождении каждой точки.

\begin{center}
\begin{tikzpicture}[x=1.6cm,y=1cm]
    \draw[->,line width=0.8pt] (-3.5,0) -- (3.8,0) node[right] {$x$};

    \fill[darkaccent] (-2,0) circle (2.4pt);
    \fill[darkaccent] (0.5,0) circle (2.4pt);
    \fill[darkaccent] (2,0) circle (2.4pt);

    \node[below=4pt] at (-2,0) {$-3$};
    \node[below=4pt] at (0.5,0) {$2$};
    \node[below=4pt] at (2,0) {$3$};

    \node[above=6pt] at (-2.8,0) {$-$};
    \node[above=6pt] at (-0.8,0) {$+$};
    \node[above=6pt] at (1.25,0) {$-$};
    \node[above=6pt] at (2.8,0) {$+$};
\end{tikzpicture}
\end{center}

Для \(\le0\) выбираем отрицательные интервалы и точки равенства:

\[
\boxed{(-\infty;-3]\cup[2;3]}.
\]

\subtopic{Алгоритм метода интервалов}

\begin{enumerate}
    \item Перенесите всё в одну сторону:
    \[
    F(x)\gtreqless0.
    \]

    \item Максимально разложите \(F(x)\) на множители.

    \item Найдите нули числителя и, если есть дробь, нули знаменателя.

    \item Отметьте критические точки на числовой прямой.

    \item Определите знак на одном удобном интервале.

    \item Расставьте остальные знаки с учётом кратностей корней.

    \item Выберите интервалы требуемого знака.

    \item Проверьте включение или исключение граничных точек.
\end{enumerate}

\subtopic{Когда знак сохраняется}

Если множитель имеет чётную степень, например

\[
(x-2)^2,
\]

при переходе через \(x=2\) знак всего произведения сохраняется.

Если степень нечётная, знак меняется.

\[
\boxed{
\begin{array}{c}
\text{нечётная кратность}\Rightarrow\text{знак меняется},\\[0.3em]
\text{чётная кратность}\Rightarrow\text{знак сохраняется}.
\end{array}
}
\]

% ============================================================
% Дробно-рациональные неравенства
% ============================================================

\topic{7. Дробно-рациональные неравенства}

Рассмотрим

\[
\frac{x^2-9}{x-1}<0.
\]

Факторизуем числитель:

\[
\frac{(x-3)(x+3)}{x-1}<0.
\]

Критические точки:

\[
-3,\qquad 1,\qquad 3.
\]

Точки \(x=-3\) и \(x=3\) обращают числитель в ноль.

Точка

\[
x=1
\]

обращает знаменатель в ноль, поэтому

\[
x=1
\]

исключается при любом знаке исходного неравенства.

\begin{center}
\begin{tikzpicture}[x=1.55cm,y=1cm]
    \draw[->,line width=0.8pt] (-3.7,0) -- (3.9,0) node[right] {$x$};

    \draw[darkaccent,line width=1pt,fill=white]
        (-2.2,0) circle (2.6pt);
    \draw[darkaccent,line width=1pt,fill=white]
        (0.4,0) circle (2.6pt);
    \draw[darkaccent,line width=1pt,fill=white]
        (2.2,0) circle (2.6pt);

    \node[below=4pt] at (-2.2,0) {$-3$};
    \node[below=4pt] at (0.4,0) {$1$};
    \node[below=4pt] at (2.2,0) {$3$};

    \node[above=6pt] at (-3,0) {$-$};
    \node[above=6pt] at (-0.9,0) {$+$};
    \node[above=6pt] at (1.3,0) {$-$};
    \node[above=6pt] at (3,0) {$+$};
\end{tikzpicture}
\end{center}

Ответ:

\[
\boxed{(-\infty;-3)\cup(1;3)}.
\]

\warning{
У дробного выражения нули знаменателя всегда являются выколотыми точками.
}

\subtopic{Правильный вид перед методом интервалов}

Желательно привести задачу к форме

\[
\frac{P(x)}{Q(x)}\gtreqless0,
\]

где слева находится одна дробь, максимально разложенная на множители.

\keyidea{
Калибровочное число подставляют уже в окончательно упрощённое выражение.
}

% ============================================================
% Текстовые задачи
% ============================================================

\topic{8. Текстовые задачи: универсальная модель}

Для задач на движение и работу удобно использовать таблицу.

\begin{table}[h]
\centering
\caption{Базовая структура}
\renewcommand{\arraystretch}{1.35}
\begin{tabularx}{0.85\linewidth}{@{}lXXX@{}}
\toprule
Случай & Величина 1 & Время & Результат \\
\midrule
1 & \(v_1\) или \(p_1\) & \(t_1\) & \(S_1\) или \(A_1\) \\
2 & \(v_2\) или \(p_2\) & \(t_2\) & \(S_2\) или \(A_2\) \\
\bottomrule
\end{tabularx}
\end{table}

Для движения:

\[
S=vt,
\qquad
t=\frac{S}{v},
\qquad
v=\frac{S}{t}.
\]

Для работы:

\[
A=pt,
\qquad
t=\frac{A}{p},
\qquad
p=\frac{A}{t},
\]

где \(p\) --- производительность.

\subtopic{Главный алгоритм}

\begin{enumerate}
    \item Определите, что требуется найти.
    \item Обозначьте искомую величину через \(x\).
    \item Заполните известные величины.
    \item Формульный столбец заполните последним.
    \item Найдите информацию из условия, которая ещё не использована.
    \item Именно эта информация обычно связывает строки таблицы уравнением.
\end{enumerate}

% ============================================================
% Велосипедист
% ============================================================

\topic{9. Движение туда и обратно}

Велосипедист проехал \(180\) км из города \(A\) в город \(B\). На обратном пути его скорость увеличилась на \(6\) км/ч, а время уменьшилось на \(1\) час.

Пусть

\[
x\ \text{км/ч}
\]

--- скорость на пути из \(A\) в \(B\).

Тогда:

\begin{table}[h]
\centering
\caption{Модель движения велосипедиста}
\renewcommand{\arraystretch}{1.4}
\begin{tabularx}{0.9\linewidth}{@{}lXXX@{}}
\toprule
Направление & Скорость, км/ч & Время, ч & Путь, км \\
\midrule
\(A\to B\) & \(x\) & \(\dfrac{180}{x}\) & \(180\) \\
\(B\to A\) & \(x+6\) & \(\dfrac{180}{x+6}\) & \(180\) \\
\bottomrule
\end{tabularx}
\end{table}

По смыслу задачи:

\[
\frac{180}{x}
-
\frac{180}{x+6}
=1.
\]

ОДЗ:

\[
x>0.
\]

\keyidea{
Если сказано, что один путь занял на \(1\) час меньше, из большего времени вычитаем меньшее.
}

% ============================================================
% Дробное время
% ============================================================

\topic{10. Работа со временем в текстовых задачах}

При вычислениях удобнее использовать обыкновенные дроби.

Например,

\[
1\text{ ч }45\text{ мин}
=
1+\frac{45}{60}
=
1+\frac34
=
\frac74\text{ ч}.
\]

\warning{
Сохраняйте точные дроби в промежуточных вычислениях. Десятичную запись удобно выполнять в самом конце, если этого требует формат ответа.
}

% ============================================================
% Поезда
% ============================================================

\topic{11. Движение протяжённого тела}

Поезд проходит мимо пешехода. Здесь движущийся объект имеет ненулевую длину.

Если объект длины \(L\) проходит мимо объекта длины \(l\), то

\[
t=\frac{L+l}{v_{\text{отн}}},
\]

где \(v_{\text{отн}}\) --- относительная скорость.

Для пешехода его собственная длина в школьной модели принимается равной нулю:

\[
l=0.
\]

\subtopic{Относительная скорость}

При движении в одном направлении:

\[
v_{\text{отн}}=|v_1-v_2|.
\]

При движении навстречу:

\[
v_{\text{отн}}=v_1+v_2.
\]

\subtopic{Пример}

Поезд движется со скоростью

\[
72\text{ км/ч},
\]

пешеход ---

\[
4{,}5\text{ км/ч},
\]

и поезд проходит мимо него за

\[
32\text{ с}.
\]

Удобно перевести скорости в метры в секунду:

\[
72\text{ км/ч}=20\text{ м/с},
\]

\[
4{,}5\text{ км/ч}=1{,}25\text{ м/с}.
\]

Относительная скорость:

\[
20-1{,}25=18{,}75\text{ м/с}.
\]

Длина поезда:

\[
L=18{,}75\cdot32=600\text{ м}.
\]

Ответ:

\[
\boxed{600\text{ м}}.
\]

\warning{
Все величины в формуле должны быть выражены в согласованных единицах.
}

% ============================================================
% Производительность
% ============================================================

\topic{12. Задачи на производительность}

Производительность --- объём работы, выполняемый за единицу времени.

Например,

\[
4\text{ л/мин}
\]

--- производительность трубы.

Связь величин:

\[
A=pt,
\]

где

\[
A\text{ --- объём работы},\qquad
p\text{ --- производительность},\qquad
t\text{ --- время}.
\]

\subtopic{Пример с трубами}

Первая труба пропускает на \(4\) л/мин меньше второй. Резервуар объёмом \(240\) л первая труба заполняет на \(5\) минут дольше.

Пусть

\[
x\text{ л/мин}
\]

--- производительность первой трубы.

Тогда производительность второй:

\[
x+4.
\]

\begin{table}[h]
\centering
\caption{Модель задачи на трубы}
\renewcommand{\arraystretch}{1.4}
\begin{tabularx}{0.9\linewidth}{@{}lXXX@{}}
\toprule
Труба & Производительность & Время & Объём \\
\midrule
Первая & \(x\) & \(\dfrac{240}{x}\) & \(240\) \\
Вторая & \(x+4\) & \(\dfrac{240}{x+4}\) & \(240\) \\
\bottomrule
\end{tabularx}
\end{table}

Так как первая труба работает дольше:

\[
\frac{240}{x}
-
\frac{240}{x+4}
=5.
\]

По смыслу:

\[
x>0.
\]

% ============================================================
% Растворы
% ============================================================

\topic{13. Растворы и смеси}

Для задач на смеси удобно использовать три величины:

\[
\boxed{
\text{масса раствора}
\quad
\text{концентрация}
\quad
\text{масса вещества}
}
\]

Формула:

\[
m_{\text{вещества}}
=
m_{\text{раствора}}\cdot c,
\]

где \(c\) записывается десятичной дробью.

Например,

\[
44\%=0{,}44.
\]

\warning{
Концентрации напрямую не складываются. Складываются массы растворов и массы растворённого вещества.
}

% ============================================================
% Два раствора
% ============================================================

\topic{14. Два раствора и система уравнений}

Имеются \(12\) кг и \(18\) кг растворов различных концентраций. После смешивания всего содержимого получается \(44\%\)-ный раствор.

Пусть

\[
x,\ y
\]

--- концентрации исходных растворов в долях единицы.

Тогда:

\begin{table}[h]
\centering
\caption{Первое смешивание}
\renewcommand{\arraystretch}{1.4}
\begin{tabularx}{0.9\linewidth}{@{}lXXX@{}}
\toprule
Раствор & Масса, кг & Концентрация & Масса вещества, кг \\
\midrule
Первый & \(12\) & \(x\) & \(12x\) \\
Второй & \(18\) & \(y\) & \(18y\) \\
Смесь & \(30\) & \(0{,}44\) & \(30\cdot0{,}44\) \\
\bottomrule
\end{tabularx}
\end{table}

Получаем:

\[
12x+18y=30\cdot0{,}44.
\]

Если равные массы двух исходных растворов дают \(50\%\)-ную смесь, обозначим каждую из этих масс через \(m\).

Тогда

\[
mx+my=2m\cdot0{,}5.
\]

Поскольку

\[
m>0,
\]

делим обе части на \(m\):

\[
x+y=1.
\]

Система:

\[
\boxed{
\begin{cases}
12x+18y=13{,}2,\\
x+y=1.
\end{cases}
}
\]

После нахождения \(x\) концентрация первого раствора в процентах равна

\[
100x\%.
\]

\keyidea{
Если в условии появляются две неизвестные концентрации, заранее ожидайте систему из двух независимых уравнений.
}

% ============================================================
% Сухофрукты
% ============================================================

\topic{15. Высушивание фруктов: ищем неизменную величину}

Свежие фрукты содержат \(84\%\) воды, высушенные --- \(36\%\) воды.

Следовательно, содержание сухого вещества составляет:

\[
100\%-84\%=16\%
\]

для свежих фруктов и

\[
100\%-36\%=64\%
\]

для высушенных.

Пусть \(x\) кг --- масса свежих фруктов, необходимая для получения \(40\) кг высушенных.

Масса сухого вещества сохраняется:

\[
0{,}16x=0{,}64\cdot40.
\]

Отсюда

\[
x=\frac{0{,}64\cdot40}{0{,}16}=160.
\]

Ответ:

\[
\boxed{160\text{ кг}}.
\]

\keyidea{
В задачах на высушивание ищите компонент, масса которого сохраняется. Обычно это сухое вещество.
}

% ============================================================
% Графики после сокращения
% ============================================================

\topic{16. График рациональной функции после сокращения}

Пусть дана функция

\[
y=
\frac{(x-4)(x-1)(x+2)}
     {(x-1)(x+2)}.
\]

Сначала фиксируем ограничения исходной функции:

\[
x\ne1,
\qquad
x\ne-2.
\]

После сокращения получаем

\[
y=x-4,
\]

при условиях

\[
x\ne1,\qquad x\ne-2.
\]

Следовательно, график --- прямая

\[
y=x-4
\]

с двумя выколотыми точками.

Их координаты:

для \(x=1\)

\[
y=1-4=-3,
\]

то есть

\[
A(1;-3);
\]

для \(x=-2\)

\[
y=-2-4=-6,
\]

то есть

\[
B(-2;-6).
\]

\begin{center}
\begin{tikzpicture}
\begin{axis}[
    width=11cm,
    height=8cm,
    axis lines=middle,
    unit vector ratio=1 1,
    xmin=-4,
    xmax=5,
    ymin=-8,
    ymax=3,
    xtick={-4,-3,...,5},
    ytick={-8,-7,...,3},
    xlabel={$x$},
    ylabel={$y$},
    samples=250,
    clip=false,
    grid=both,
    major grid style={line width=0.25pt},
    minor grid style={line width=0.15pt},
]
    \addplot[
        domain=-4:5,
        accent,
        line width=1.3pt
    ] {x-4};

    \addplot[
        only marks,
        mark=o,
        mark size=3pt,
        line width=1pt,
        darkaccent
    ]
    coordinates {(1,-3) (-2,-6)};

    \node[anchor=south west,text=darkaccent]
        at (axis cs:1,-3) {$A(1;-3)$};

    \node[anchor=north east,text=darkaccent]
        at (axis cs:-2,-6) {$B(-2;-6)$};
\end{axis}
\end{tikzpicture}
\end{center}

\warning{
После сокращения запрещённые значения \(x\) сохраняются. На графике им соответствуют выколотые точки.
}

% ============================================================
% y = m
% ============================================================

\topic{17. Параметр \(m\) и горизонтальная прямая}

Уравнение

\[
y=m
\]

задаёт горизонтальную прямую.

Для графика

\[
y=x-4,
\qquad
x\ne1,\ -2,
\]

горизонтальная прямая не имеет с графиком общих точек только тогда, когда проходит через уровень одной из выколотых точек.

У первой выколотой точки:

\[
y=-3.
\]

У второй:

\[
y=-6.
\]

Следовательно,

\[
\boxed{m=-6;\ -3}.
\]

\subtopic{Различайте две семьи прямых}

\[
y=kx
\]

--- прямая, проходящая через начало координат.

\[
y=m
\]

--- горизонтальная прямая.

\keyidea{
Если требуется найти значения \(m\) в уравнении \(y=m\), ищите особые значения координаты \(y\).
}

% ============================================================
% Быстрый выбор метода
% ============================================================

\topic{18. Быстрая схема выбора метода}

\begin{center}
\begin{tikzpicture}[
    every node/.style={font=\small},
    block/.style={
        draw=darkaccent,
        rounded corners=2pt,
        inner sep=7pt,
        text width=11.3cm,
        align=left
    }
]

\node[block] (a) at (0,0)
{\textbf{4 слагаемых, похожие пары}
\(\;\Rightarrow\;\)
проверить группировку и двойное вынесение.};

\node[block] (b) at (0,-1.75)
{\textbf{Скобка в квадрате и число}
\(\;\Rightarrow\;\)
изолировать квадрат:
\(A(x)^2=c\).};

\node[block] (c) at (0,-3.5)
{\textbf{Повторяющийся сложный элемент}
\(\;\Rightarrow\;\)
проверить замену переменной.};

\node[block] (d) at (0,-5.25)
{\textbf{Произведение равно нулю в системе}
\(\;\Rightarrow\;\)
разбить на случаи и выполнить подстановку.};

\node[block] (e) at (0,-7)
{\textbf{Неравенство}
\(\;\Rightarrow\;\)
перенести всё в одну сторону, факторизовать, применить метод интервалов.};

\node[block] (f) at (0,-8.75)
{\textbf{Текстовая задача}
\(\;\Rightarrow\;\)
ввести \(x\), построить таблицу, заполнить формульный столбец, использовать оставшуюся связь из условия.};

\node[block] (g) at (0,-10.5)
{\textbf{Рациональная функция}
\(\;\Rightarrow\;\)
сначала ограничения, затем сокращение, построение упрощённого графика и выколотых точек.};

\end{tikzpicture}
\end{center}

% ============================================================
% Типичные ошибки
% ============================================================

\topic{19. Что обязательно проверять}

\begin{enumerate}
    \item \textbf{Знаменатели.}
    \[
    Q(x)\ne0.
    \]

    \item \textbf{Сокращение дробей.}
    Исключённые значения исходного знаменателя сохраняются.

    \item \textbf{Квадраты.}
    \[
    A(x)^2\ge0.
    \]

    \item \textbf{Замена переменной.}
    Проверяйте ограничения на новую переменную, если они возникают.

    \item \textbf{Системы.}
    Ответ записывается упорядоченными парами.

    \item \textbf{Метод интервалов.}
    Нули знаменателя всегда исключаются.

    \item \textbf{Кратность корней.}
    При чётной кратности знак сохраняется.

    \item \textbf{Текстовые задачи.}
    Проверяйте физический и прикладной смысл найденного значения:
    \[
    v>0,\qquad t>0,\qquad p>0.
    \]

    \item \textbf{Концентрации.}
    В формулах используйте доли единицы:
    \[
    44\%=0{,}44.
    \]

    \item \textbf{Графики.}
    Указывайте масштаб, подписи осей и точные координаты выколотых точек.
\end{enumerate}

% ============================================================
% Итоговый чек-лист
% ============================================================

\topic{20. Итоговый чек-лист перед экзаменом}

Перед сдачей решения проверьте:

\begin{itemize}
    \item выражение максимально упрощено;
    \item все общие множители вынесены;
    \item ОДЗ записана;
    \item запрещённые значения нигде не потеряны;
    \item уравнение со скобкой в квадрате решено коротким способом, если структура это позволяет;
    \item при замене выполнен обратный переход;
    \item найденные корни проверены на ограничения;
    \item система завершена упорядоченными парами;
    \item в неравенстве правильно отмечены закрашенные и выколотые точки;
    \item ось подписана;
    \item единицы измерения согласованы;
    \item для растворов проценты переведены в доли;
    \item в задачах на высушивание найдена сохраняющаяся величина;
    \item у рационального графика нанесены все выколотые точки;
    \item записана строка \textbf{Ответ:}.
\end{itemize}

% ============================================================
% ТРЕНИРОВОЧНЫЙ БЛОК
% ============================================================

\newpage

\thispagestyle{fancy}

\begin{center}
{\LARGE\bfseries\textcolor{darkaccent}{Тренировочный блок}}\\[0.3em]
{\large Путь А}
\end{center}

{\color{accent}\rule{\linewidth}{0.8pt}}

\vspace{0.6em}

Решите задачи без подсказок. Первые три задания имеют средний уровень, следующие шесть --- уровень выше среднего, последнее задание --- сложное комплексное.

\begin{enumerate}

% 1
\item
Разложите на множители:

\[
x^2+3x-4x-12.
\]

% 2
\item
Решите уравнение:

\[
(x-7)^2-13=0.
\]

% 3
\item
Решите неравенство:

\[
(x+4)(x-1)(x-5)\le0.
\]

% 4
\item
Решите уравнение:

\[
\frac{1}{(x+2)^2}
-\frac{7}{x+2}
+10=0.
\]

% 5
\item
Решите систему:

\[
\begin{cases}
(x-2)(y+1)=0,\\[0.4em]
\dfrac{x+y}{x+y-1}=2.
\end{cases}
\]

% 6
\item
Решите неравенство:

\[
\frac{(x-4)(x+2)}{x-1}\ge0.
\]

% 7
\item
Автомобиль проехал \(240\) км из города \(A\) в город \(B\).
На обратном пути скорость была на \(20\) км/ч больше, поэтому дорога заняла на \(1\) час меньше.
Найдите скорость автомобиля на пути из \(A\) в \(B\).

% 8
\item
Первая труба пропускает на \(6\) л/мин меньше второй.
Резервуар объёмом \(360\) л первая труба заполняет на \(5\) минут дольше второй.
Найдите производительность первой трубы.

% 9
\item
Свежие ягоды содержат \(88\%\) воды, высушенные --- \(40\%\) воды.
Сколько килограммов свежих ягод требуется для получения \(18\) кг высушенных?

% 10
\item
Дана функция

\[
y=
\frac{(x-5)(x-2)(x+3)}
     {(x-2)(x+3)}.
\]

\begin{enumerate}[label=\alph*)]
    \item Постройте её график с учётом всех ограничений.
    \item Найдите координаты всех выколотых точек.
    \item Найдите все значения \(m\), при которых прямая
    \[
    y=m
    \]
    не имеет с графиком функции общих точек.
\end{enumerate}

\end{enumerate}

% ============================================================
% ОТВЕТЫ
% ============================================================

\newpage

\begin{center}
{\LARGE\bfseries\textcolor{darkaccent}{Ответы к тренировочному блоку}}
\end{center}

{\color{accent}\rule{\linewidth}{0.8pt}}

\vspace{0.7em}

\renewcommand{\arraystretch}{1.55}

\begin{longtable}{@{}p{1.2cm}p{0.79\linewidth}@{}}
\toprule
\textbf{№} & \textbf{Ответ} \\
\midrule
\endfirsthead

\toprule
\textbf{№} & \textbf{Ответ} \\
\midrule
\endhead

1 &
\[
(x+3)(x-4).
\]
\\

2 &
\[
x=7-\sqrt{13},
\qquad
x=7+\sqrt{13}.
\]
\\

3 &
\[
(-\infty;-4]\cup[1;5].
\]
\\

4 &
\[
x=-\frac32,
\qquad
x=-\frac{19}{10}.
\]
\\

5 &
\[
(2;0),
\qquad
(-2;-1).
\]
\\

6 &
\[
[-2;1)\cup[4;+\infty).
\]
\\

7 &
\[
60\text{ км/ч}.
\]
\\

8 &
\[
18\text{ л/мин}.
\]
\\

9 &
\[
90\text{ кг}.
\]
\\

10 &
После сокращения:
\[
y=x-5,
\qquad
x\ne2,\quad x\ne-3.
\]

Выколотые точки:
\[
A(2;-3),
\qquad
B(-3;-8).
\]

Значения параметра:
\[
m=-8,\qquad m=-3.
\]
\\

\bottomrule
\end{longtable}

\vfill

\begin{center}
\textcolor{softgray}{
\small
Главная цель повторения --- довести распознавание типовой структуры задачи до автоматизма.
}
\end{center}

\end{document}